\section*{Chapitre \ref{ch:b3:many-electron-atoms} --- Atomes polyélectroniques et symboles de termes}

\begin{solution}{exo:b3:many-electron-atoms:1}
\[
\Psi = \frac{1}{\sqrt6}\begin{vmatrix}1s\alpha(1) & 1s\beta(1) & 2s\alpha(1)\\
1s\alpha(2) & 1s\beta(2) & 2s\alpha(2)\\ 1s\alpha(3) & 1s\beta(3) & 2s\alpha(3)
\end{vmatrix}.
\]
$1s$ n'offre que deux \omterm{def:b3:many-electron-atoms:spin-orbital}{spin-orbitales}, $1s\alpha$ et $1s\beta$~; un troisième
électron en répéterait une, ce qui rendrait deux colonnes égales~: le
déterminant est nul.
\end{solution}

\begin{solution}{exo:b3:many-electron-atoms:2}
$\binom62 = 15$, $\binom63 = 20$, $\binom{10}{2} = 45$. $p^3$~: ${}^4$S (4) +
${}^2$D (10) + ${}^2$P (6) = 20. $d^2$~: ${}^3$F (21) + ${}^3$P (9) + ${}^1$G (9) +
${}^1$D (5) + ${}^1$S (1) = 45.
\end{solution}

\begin{solution}{exo:b3:many-electron-atoms:3}
N ($2p^3$, à moitié pleine)~: \termsym{4}{S}{3/2}. O ($2p^4$, plus qu'à
moitié pleine)~:
\termsym{3}{P}{2}. F ($2p^5$)~: \termsym{2}{P}{3/2}. \ce{Ti^2+} ($3d^2$)~:
\termsym{3}{F}{2}. \ce{Fe^3+} ($3d^5$, à moitié pleine, tous les spins
parallèles, $L = 0$)~:
\termsym{6}{S}{5/2}.
\end{solution}

\begin{solution}{exo:b3:many-electron-atoms:4}
${}^2$D~: $J = \frac52$ (6 états) et $\frac32$ (4)~; $6 + 4 = 10 = 5 \times 2$.
${}^3$P~: $J = 2$ (5), 1 (3), 0 (1)~; $5 + 3 + 1 = 9 = 3 \times 3$.
\end{solution}

\begin{solution}{exo:b3:many-electron-atoms:5}
Intervalles $16.4$ ($J = 1 - 0$) et $27.0$ ($2 - 1$)~: rapport 1.65 au lieu
de 2.
$A = 16.4/1 = \qty{16.4}{cm^{-1}}$ d'après le premier, $27.0/2 = \qty{13.5}{cm^{-1}}$
d'après le second. Une constante $A$ unique ne convient pas~: les niveaux
sont légèrement mélangés avec ceux d'autres termes (${}^1$D, ${}^1$S), et
le \omterm{def:b3:many-electron-atoms:russell-saunders}{couplage de Russell--Saunders} est une approximation.
\end{solution}

\begin{solution}{exo:b3:many-electron-atoms:6}
$\lambda = 10^7/\tilde\nu$~: \qty{589.76}{nm} (D$_1$) et \qty{589.16}{nm}
(D$_2$). Écart $\qty{17.20}{cm^{-1}} = \qty{2.13}{meV}$. $E(\frac32) -
E(\frac12) = \frac32A$, donc $A = \qty{11.47}{cm^{-1}}$.
\end{solution}

\begin{solution}{exo:b3:many-electron-atoms:7}
$3p \to 3s$~: permise ($\Delta l = -1$). $3d \to 3s$~: interdite ($\Delta l = -2$).
$3d \to 3p$~: permise. $4s \to 3s$~: interdite ($\Delta l = 0$, pas de
changement de parité). $4s \to 3p$~: permise.
\end{solution}

\begin{solution}{exo:b3:many-electron-atoms:8}
Moyenne de ${}^3$P~: $(5 \times 169086.76 + 3 \times 169086.84 + 169087.83)/9 =
\qty{169086.91}{cm^{-1}}$. Écart à ${}^1$P~: \qty{2047.98}{cm^{-1}} $=
\qty{0.254}{eV}$, donc $K(1s,2p) = \qty{0.127}{eV}$, trois fois moins que
$K(1s,2s)$. L'\omterm{def:b3:many-electron-atoms:exchange}{intégrale d'échange} mesure le recouvrement des densités des
deux orbitales~; l'orbitale $2s$ pénètre dans la région de $1s$ (elle a
de la densité près du noyau), l'orbitale $2p$ s'annule au noyau et
recouvre moins.
\end{solution}

\begin{solution}{exo:b3:many-electron-atoms:9}
Le plus grand $M_L = 4$ exige les deux électrons dans $m_l = 2$, de spins
opposés~: un singulet, ${}^1$G (9 états). Le suivant, $M_L = 3$, provient
de $m_l = 2, 1$, avec
$M_S = 1$ possible~: ${}^3$F (21). Il reste $M_L = 2$ avec $M_S = 0$
seulement~:
${}^1$D (5). Il reste $M_L = 1$ avec $M_S = 1$~: ${}^3$P (9). Il reste un
état avec
$M_L = M_S = 0$~: ${}^1$S. Total 45.
\end{solution}

\begin{solution}{exo:b3:many-electron-atoms:10}
En $\mathbf r_1 = \mathbf r_2 = \mathbf r$, la fonction triplet vaut
$\frac{1}{\sqrt2}[1s(\mathbf r)2s(\mathbf r) - 2s(\mathbf r)1s(\mathbf r)] = 0$~; la
fonction singulet vaut $\sqrt2\,1s(\mathbf r)2s(\mathbf r) \ne 0$. Dans
le triplet, les électrons ne se trouvent jamais au même point et sont en
moyenne plus éloignés, donc leur répulsion est plus faible~:
$E_+ - E_- = 2K > 0$.
\end{solution}

\begin{solution}{exo:b3:many-electron-atoms:11}
$d^8$ est plus qu'à moitié pleine~: la structure est inversée,
$J = L + S = 4$ est le plus bas. Intervalles~: $E(3) - E(4) = 1360.7$ et
$E(2) - E(3) = 908.9$~; rapport
1.50, contre $4/3 = 1.33$ selon Landé. $|A| = 1360.7/4 = \qty{340}{cm^{-1}}$
($A < 0$). Le \omterm{def:b3:many-electron-atoms:spin-orbit}{couplage spin--orbite} est bien plus fort dans cet ion plus
lourd que dans le carbone (\qty{16}{cm^{-1}}).
\end{solution}

\begin{solution}{exo:b3:many-electron-atoms:12}
$\sum_J(2J+1)[J(J+1) - L(L+1) - S(S+1)] = 0$ (on peut dénombrer les
$(2L+1)(2S+1)$ états dans l'une ou l'autre base, et $\sum M_LM_S = 0$ sur
les états découplés).
Pour ${}^3$P~: $1(0 - 4) + 3(2 - 4) + 5(6 - 4) = -4 - 6 + 10 = 0$. La
moyenne pondérée du carbone vaut $(0 + 3 \times 16.42 + 5 \times 43.41)/9 = \qty{29.6}{cm^{-1}}$~:
l'énergie qu'aurait le terme sans \omterm{def:b3:many-electron-atoms:spin-orbit}{couplage spin--orbite}.
\end{solution}

\begin{solution}{pb:b3:many-electron-atoms:1}
\textbf{1.} $1s\alpha$, $1s\beta$, $2s\alpha$, $2s\beta$.
\textbf{2.} $|1s\alpha\,2s\alpha|$, $|1s\alpha\,2s\beta|$, $|1s\beta\,2s\alpha|$,
$|1s\beta\,2s\beta|$.
\textbf{3.} $|1s\alpha\,2s\alpha| = \frac{1}{\sqrt2}[1s(1)2s(2) - 2s(1)1s(2)]\,
\alpha(1)\alpha(2)$, en développant le déterminant $2\times2$.
\textbf{4.} Leur somme vaut $\frac{1}{\sqrt2}[1s(1)2s(2) - 2s(1)1s(2)]\cdot
\frac{1}{\sqrt2}[\alpha(1)\beta(2) + \beta(1)\alpha(2)]$ (multipliée par
$\sqrt2$, puis normée)~; leur différence donne $\frac{1}{\sqrt2}[1s(1)2s(2) +
2s(1)1s(2)]\cdot\frac{1}{\sqrt2}[\alpha(1)\beta(2) - \beta(1)\alpha(2)]$.
\textbf{5.} Fonction d'espace symétrique avec la fonction de spin
antisymétrique ($S = 0$)~;
fonction d'espace antisymétrique avec les trois fonctions de spin
symétriques ($S = 1$).
\textbf{6.} $2S + 1 = 1$ et 3~: un et trois états de même énergie.
\textbf{7.} $[\ce{Ar}]\,3d^2$~; $\binom{10}{2} = 45$ déterminants.
\textbf{8.} Valeur maximale $S = 1$, valeur maximale $L = 2 + 1 = 3$,
sous-couche moins qu'à moitié pleine~:
\termsym{3}{F}{2}, le niveau à 0 dans les données.
\textbf{9.} $21 + 9 + 9 + 5 + 1 = 45$.
\textbf{10.} Normale ($J = 2$ le plus bas). Intervalles 184.9 et 235.5~:
rapport 1.27
contre $4/3 = 1.33$~; la règle tient à 5\,\% près.
\textbf{11.} ${}^3$F~: $(5 \times 0 + 7 \times 184.9 + 9 \times 420.4)/21 =
\qty{241.8}{cm^{-1}}$~; ${}^3$P~: $(10538.4 + 3 \times 10603.6 + 5 \times
10721.2)/9 = \qty{10661.7}{cm^{-1}}$.
\textbf{12.} $15B = \qty{10419.9}{cm^{-1}}$, $B = \qty{695}{cm^{-1}}$.
\textbf{13.} $\Delta S = 0$.
\textbf{14.} $\Delta S = 1$ est interdit, et $1s2s \to 1s^2$ a $\Delta l = 0$
(pas de changement de parité~; ${}^3$S$_1 \to {}^1$S$_0$ a aussi
$\Delta L = 0$ entre deux états S). L'état est métastable~: il vit bien
plus longtemps que les états excités ordinaires.
\textbf{15.} $169086.91 - 159855.97 = \qty{9230.94}{cm^{-1}}$~:
\qty{1083.3}{nm}, dans le proche infrarouge.
\textbf{16.} $171134.89 - 166277.44 = \qty{4857.45}{cm^{-1}}$~: \qty{2058.7}{nm}.
\textbf{17.} $\qty{171134.89}{cm^{-1}}$~: \qty{58.43}{nm}, dans
l'ultraviolet du vide (absorbé par l'air).
\textbf{18.} Chaque famille a ses propres raies et son propre état le plus
bas, et les deux ne communiquent jamais par la lumière~: elles se
comportent comme deux gaz aux deux spectres distincts.
\textbf{19.} $E_\pm = E_{1s} + E_{2s} + J \pm K$ (singulet $+$, triplet $-$).
\textbf{20.} Le triplet~: sa fonction spatiale s'annule quand les
électrons se rencontrent, si bien qu'ils se repoussent moins.
\textbf{21.} \qty{6421.47}{cm^{-1}}, $6421.47/8065.54 = \qty{0.796}{eV}$.
\textbf{22.} Écart $\qty{2047.98}{cm^{-1}} = \qty{0.254}{eV}$, donc $K(1s,2p) =
\qty{0.127}{eV}$.
\textbf{23.} $2s$ pénètre dans la région de $1s$ et la recouvre davantage
que $2p$.
\textbf{24.} Oui~: dans les deux configurations, le triplet ($S$ plus
grand) est le plus bas.
\textbf{25.} $K = \frac12 \times \qty{0.796}{eV}$~: \textbf{l'\omterm{def:b3:many-electron-atoms:exchange}{intégrale
d'échange} $K(1s,2s)$ de l'hélium vaut \qty{0.398}{eV}}, mesurée comme la
moitié de l'écart singulet--triplet.
\end{solution}
